% GENERATED FILE. Edit canonical lessons, then run npm run generate:books.
% canonical-source-hash: fnv1a64:a43ea869413bed27
% canonical-lessons: GU-C240-chod, GU-C240-bij, GU-C240-ropvu, GU-C240-sicvu, GU-C240-gulab

\chapter{In the garden}
\label{ch:gu-in-the-garden}
\hlchaptermodality{240}

\begin{chapteropening}
\textbf{By the end of this chapter:} \emph{I can say what I planted and watered.}

Complete the last lesson of chapter 240: I can say what I planted and watered.
\end{chapteropening}

\section[\texorpdfstring{\gu{છોડ} (choḍ)}{છોડ (choḍ)}]{\gu{છોડ} (choḍ) — plant, shrub}

\label{lesson:GU-C240-chod}



\begin{quote}
Before the new one: say the Gujarati for to sting, then the Gujarati for to slip away.
\end{quote}

\subsection*{You'll want to know: \gu{છોડ}}
\textbf{\gu{છોડ}} — \emph{choḍ} — \textquotedblleft{}plant, shrub\textquotedblright{}.

\begin{quote}
\textbf{\gu{મેં} \gu{છોડને} \gu{પાણી} \gu{આપ્યું}.} — \emph{mẽ choḍne pāṇī āpyũ} — \textquotedblleft{}I gave the plant water.\textquotedblright{}
\end{quote}

A masculine noun.

\begin{cousinweb}
The origin of \textbf{\gu{છોડ}} (\emph{choḍ}), \textquotedblleft{}plant\textquotedblright{}, is not settled: CDIAL lists no Gujarati form like it among its words for plants.
\end{cousinweb}

\begin{grammarlens}[title={What you've built}]
One more word for everyday life at A2.
\end{grammarlens}

\subsection*{Your turn}
\begin{itemize}
\raggedright
  \item \emph{Say it:} \emph{choḍ}
  \item \emph{Say it:} \emph{choḍ}, once more
  \item \emph{Say it:} \emph{mẽ choḍne pāṇī āpyũ} — I gave the plant water.
  \item \emph{Recall:} say the Gujarati for to sting, then the Gujarati for to slip away, then say \emph{choḍ} again
\end{itemize}

\begin{tcolorbox}[breakable,colback=teal!4,colframe=teal!35!black,title={Before you move on}]
What does \gu{છોડ} mean? (\textquotedblleft{}Plant, shrub\textquotedblright{}.) Say it once more.
\end{tcolorbox}
\section[\texorpdfstring{\gu{બીજ} (bīj)}{બીજ (bīj)}]{\gu{બીજ} (bīj) — seed}

\label{lesson:GU-C240-bij}



\begin{quote}
Before the new one: say the Gujarati for to slip away, then the Gujarati for plant.
\end{quote}

\subsection*{You'll want to know: \gu{બીજ}}
\textbf{\gu{બીજ}} — \emph{bīj} — \textquotedblleft{}seed\textquotedblright{}.

\begin{quote}
\textbf{\gu{મેં} \gu{ખેતરમાં} \gu{બીજ} \gu{વાવ્યાં}.} — \emph{mẽ khetarmā̃ bīj vāvyā̃} — \textquotedblleft{}I sowed seeds in the field.\textquotedblright{}
\end{quote}

Not \textbf{\gu{બીજું}}, second.

\begin{cousinweb}
\textbf{\gu{બીજ}} (\emph{bīj}) is Sanskrit \emph{bīja}, \textquotedblleft{}seed\textquotedblright{}, taken over almost whole, a tatsama. The inherited forms of the same word are Gujarati \emph{bī} and \emph{bīyũ}, \textquotedblleft{}seed\textquotedblright{} (CDIAL 9250).
\end{cousinweb}

\begin{grammarlens}[title={What you've built}]
One more word for everyday life at A2.
\end{grammarlens}

\subsection*{Your turn}
\begin{itemize}
\raggedright
  \item \emph{Say it:} \emph{bīj}
  \item \emph{Say it:} \emph{bīj}, once more
  \item \emph{Say it:} \emph{mẽ khetarmā̃ bīj vāvyā̃} — I sowed seeds in the field.
  \item \emph{Recall:} say the Gujarati for to slip away, then the Gujarati for plant, then say \emph{bīj} again
\end{itemize}

\begin{tcolorbox}[breakable,colback=teal!4,colframe=teal!35!black,title={Before you move on}]
What does \gu{બીજ} mean? (\textquotedblleft{}Seed\textquotedblright{}.) Say it once more.
\end{tcolorbox}
\section[\texorpdfstring{\gu{રોપવું} (ropvũ)}{રોપવું (ropvũ)}]{\gu{રોપવું} (ropvũ) — to plant, to set out}

\label{lesson:GU-C240-ropvu}



\begin{quote}
Before the new one: say the Gujarati for plant, then the Gujarati for seed.
\end{quote}

\subsection*{You'll want to know: \gu{રોપવું}}
\textbf{\gu{રોપવું}} — \emph{ropvũ} — \textquotedblleft{}to plant, to set out\textquotedblright{}.

\begin{quote}
\textbf{\gu{તેણે} \gu{બગીચામાં} \gu{છોડ} \gu{રોપ્યો}.} — \emph{teṇe bagīcāmā̃ choḍ ropyo} — \textquotedblleft{}She planted a plant in the garden.\textquotedblright{}
\end{quote}

\textbf{\gu{વાવવું}} is to sow seed; \textbf{\gu{રોપવું}} is to put a young plant in the ground.

\begin{cousinweb}
\textbf{\gu{રોપવું}} (\emph{ropvũ}) is a tadbhava. CDIAL lists it, glossed \textquotedblleft{}to plant\textquotedblright{}, under \emph{ropyate}, \textquotedblleft{}is planted\textquotedblright{}, the passive of Sanskrit \emph{ropayati}, \textquotedblleft{}plants\textquotedblright{}, in its entry for \emph{*rupati}, \textquotedblleft{}pierces, makes a hole, plants\textquotedblright{} (CDIAL 10783).
\end{cousinweb}

\begin{grammarlens}[title={What you've built}]
One more word for everyday life at A2.
\end{grammarlens}

\subsection*{Your turn}
\begin{itemize}
\raggedright
  \item \emph{Say it:} \emph{ropvũ}
  \item \emph{Say it:} \emph{ropvũ}, once more
  \item \emph{Say it:} \emph{teṇe bagīcāmā̃ choḍ ropyo} — She planted a plant in the garden.
  \item \emph{Recall:} say the Gujarati for plant, then the Gujarati for seed, then say \emph{ropvũ} again
\end{itemize}

\begin{tcolorbox}[breakable,colback=teal!4,colframe=teal!35!black,title={Before you move on}]
What does \gu{રોપવું} mean? (\textquotedblleft{}To plant, to set out\textquotedblright{}.) Say it once more.
\end{tcolorbox}
\section[\texorpdfstring{\gu{સીંચવું} (sī̃cvũ)}{સીંચવું (sī̃cvũ)}]{\gu{સીંચવું} (sī̃cvũ) — to irrigate, to water}

\label{lesson:GU-C240-sicvu}



\begin{quote}
Before the new one: say the Gujarati for seed, then the Gujarati for to plant.
\end{quote}

\subsection*{You'll want to know: \gu{સીંચવું}}
\textbf{\gu{સીંચવું}} — \emph{sī̃cvũ} — \textquotedblleft{}to irrigate, to water\textquotedblright{}.

\begin{quote}
\textbf{\gu{મેં} \gu{સાંજે} \gu{છોડ} \gu{સીંચ્યો}.} — \emph{mẽ sā̃je choḍ sī̃cyo} — \textquotedblleft{}I watered the plant in the evening.\textquotedblright{}
\end{quote}

\begin{cousinweb}
\textbf{\gu{સીંચવું}} (\emph{sī̃cvũ}) is a tadbhava. It comes from Sanskrit \emph{siñcati}, \textquotedblleft{}pours out, sprinkles\textquotedblright{}, and CDIAL glosses the Gujarati verb \textquotedblleft{}to irrigate\textquotedblright{} (CDIAL 13394).
\end{cousinweb}

\begin{grammarlens}[title={What you've built}]
One more word for everyday life at A2.
\end{grammarlens}

\subsection*{Your turn}
\begin{itemize}
\raggedright
  \item \emph{Say it:} \emph{sī̃cvũ}
  \item \emph{Say it:} \emph{sī̃cvũ}, once more
  \item \emph{Say it:} \emph{mẽ sā̃je choḍ sī̃cyo} — I watered the plant in the evening.
  \item \emph{Recall:} say the Gujarati for seed, then the Gujarati for to plant, then say \emph{sī̃cvũ} again
\end{itemize}

\begin{tcolorbox}[breakable,colback=teal!4,colframe=teal!35!black,title={Before you move on}]
What does \gu{સીંચવું} mean? (\textquotedblleft{}To irrigate, to water\textquotedblright{}.) Say it once more.
\end{tcolorbox}
\section[\texorpdfstring{\gu{ગુલાબ} (gulāb)}{ગુલાબ (gulāb)}]{\gu{ગુલાબ} (gulāb) — rose}

\label{lesson:GU-C240-gulab}



\begin{quote}
Before the new one: say the Gujarati for to plant, then the Gujarati for to irrigate.
\end{quote}

\subsection*{You'll want to know: \gu{ગુલાબ}}
\textbf{\gu{ગુલાબ}} — \emph{gulāb} — \textquotedblleft{}rose\textquotedblright{}.

\begin{quote}
\textbf{\gu{મેં} \gu{મારી} \gu{માતાને} \gu{ગુલાબ} \gu{આપ્યું}.} — \emph{mẽ mārī mātāne gulāb āpyũ} — \textquotedblleft{}I gave my mother a rose.\textquotedblright{}
\end{quote}

\begin{cousinweb}
\textbf{\gu{ગુલાબ}} (\emph{gulāb}) is a loan, usually traced to Persian \emph{gulāb}, \textquotedblleft{}rose water\textquotedblright{}, built from \emph{gul}, \textquotedblleft{}rose\textquotedblright{}, and \emph{āb}, \textquotedblleft{}water\textquotedblright{}.
\end{cousinweb}

\begin{grammarlens}[title={What you've built}]
One more word for everyday life at A2.
\end{grammarlens}

\subsection*{Your turn}
\begin{itemize}
\raggedright
  \item \emph{Say it:} \emph{gulāb}
  \item \emph{Say it:} \emph{gulāb}, once more
  \item \emph{Say it:} \emph{mẽ mārī mātāne gulāb āpyũ} — I gave my mother a rose.
  \item \emph{Recall:} say the Gujarati for to plant, then the Gujarati for to irrigate, then say \emph{gulāb} again
\end{itemize}

\begin{tcolorbox}[breakable,colback=teal!4,colframe=teal!35!black,title={Before you move on}]
What does \gu{ગુલાબ} mean? (\textquotedblleft{}Rose\textquotedblright{}.) Say it once more.
\end{tcolorbox}
